\documentclass[11pt,a4paper]{article}

% =============================================================
% PACOTES
% =============================================================
\usepackage[T1]{fontenc}
\usepackage[utf8]{inputenc}
\usepackage[brazil]{babel}
\usepackage{lmodern}
\usepackage{microtype}
\usepackage[a4paper, margin=2.3cm]{geometry}
\usepackage{amsmath, amssymb}
\usepackage{booktabs}
\usepackage{array}
\usepackage{graphicx}
\usepackage{xcolor}
\usepackage{enumitem}
\usepackage{parskip}       % parágrafos sem recuo, separados por um espaço
\usepackage{needspace}     % evita título sozinho no pé da página
\usepackage{titlesec}
\usepackage{tcolorbox}
\tcbuselibrary{breakable, skins}
\usepackage[hidelinks]{hyperref}

% seção só começa se couber um pedaço do texto na mesma página
\newcommand{\sectionbreak}{\needspace{14\baselineskip}}
\newcommand{\subsectionbreak}{\needspace{8\baselineskip}}

% coluna de tabela com texto alinhado à esquerda
\newcolumntype{L}[1]{>{\raggedright\arraybackslash}p{#1}}

% =============================================================
% CORES E CAIXAS: cinza e um único azul-escuro
% =============================================================
\definecolor{Destaque}{RGB}{16, 37, 66}
\definecolor{Fundo}{gray}{0.965}

\tcbset{guia/.style={enhanced, breakable, sharp corners, frame hidden, colback=Fundo,
                    borderline west={1.2pt}{0pt}{black!35},
                    left=3mm, right=2mm, top=1.5mm, bottom=1.5mm, before skip=6pt, after skip=8pt}}
\newtcolorbox{exemplo}[1]{guia, before upper={{\small\bfseries Exemplo: #1}\par}}
\newtcolorbox{defesa}{guia, no borderline, borderline west={1.6pt}{0pt}{Destaque},
                      before upper={{\small\bfseries Na defesa}\par}}

% título interno, numa linha própria
\newcommand{\bloco}[1]{\par\needspace{4\baselineskip}{\bfseries #1}\par\nopagebreak}

% destaque de divergência entre políticas (modo matemático)
\newcommand{\dif}[1]{\boldsymbol{#1}}

% figura opcional: só entra se o arquivo existir na pasta imagens/
\newcommand{\figura}[3]{%
  \IfFileExists{imagens/#1}{%
    \begin{figure}[htbp]
      \centering
      \includegraphics[width=#2]{imagens/#1}
      \caption{#3}
    \end{figure}%
  }{}%
}

\setlist[itemize]{leftmargin=*, itemsep=2pt}
\setlist[enumerate]{leftmargin=*, itemsep=2pt}

% =============================================================
\title{\textbf{Guia de Estudos}\\[4pt]
\large Aprendizado por Reforço Aplicado ao \textit{Trading} de BTC-USD}
\author{Tiago Macedo, Viktor Mayer e Enzo Guedes\\
\small IBMEC --- Professor: Rodrigo Luna}
\date{Novembro de 2026}

\begin{document}

\maketitle

\textbf{Como usar este guia.} Os assuntos seguem a ordem da apresentação. Cada seção traz a ideia, a fórmula, os números do projeto e, quando ajuda, um exemplo e o que responder se o professor perguntar. Todos os números vêm do notebook \texttt{RL\_TRADE.ipynb}. A última seção reúne as perguntas mais prováveis, com respostas curtas.

\tableofcontents

% =============================================================
\section*{Números-chave}
\addcontentsline{toc}{section}{Números-chave}

\begin{center}
\small
\begin{tabular}{l L{10.6cm}}
  \toprule
  \textbf{Item} & \textbf{Valor} \\
  \midrule
  Ativo e frequência        & BTC-USD, diário (Yahoo Finance), fechamento às 21h de Brasília (0h UTC) \\
  Treino                    & 06/11/2014 a 31/12/2021 \\
  Teste                     & 01/01/2022 a 18/11/2026 (1.783 dias) \\
  Corretagem $\kappa$       & 0,1\% por unidade trocada (inverter a posição paga 0,2\%) \\
  Fator de desconto         & $\gamma = 0{,}8$ (a recompensa do 5º dia pesa 32,8\%) \\
  Limiares (tercis)         & $\tau_1 = -0{,}0234$ \quad $\tau_2 = 0{,}0953$ \\
  Espaço discreto           & 27 estados, 3 ações, 81 pares estado-ação \\
  \textit{Value Iteration}  & 98 iterações (tolerância $< 10^{-12}$) \\
  Monte Carlo e Q-learning  & 30.000 episódios de 84 dias, $\varepsilon$ de 1 a 0,05, $H = 42$ \\
  Sementes                  & 42 (principal); 42 a 46 na análise de variabilidade \\
  Taxa livre de risco       & T-bill de 3 meses: 0,78\% a.a. no treino, 3,94\% a.a. no teste \\
  Sharpe no teste           & VI 0,46 \quad MC 0,44 \quad QL 0,26 \quad \textit{Buy \& Hold} 0,04 \\
  \bottomrule
\end{tabular}
\end{center}

\section*{Símbolos}
\addcontentsline{toc}{section}{Símbolos}

\begin{center}
\small
\begin{tabular}{l L{12.5cm}}
  \toprule
  \textbf{Símbolo} & \textbf{Significado} \\
  \midrule
  $c_t$                     & preço de fechamento no dia $t$ \\
  $r_{t+1}$                 & retorno simples do dia seguinte, $c_{t+1}/c_t - 1$ \\
  $\rho_t$                  & log-retorno do dia, $\ln(c_t/c_{t-1})$ (usado pelos modelos econométricos) \\
  $x_t$                     & sinal do estado: $\ln(c_t/\mathrm{MA}_{50})$, distância do preço à média móvel de 50 dias \\
  $\phi(x)$                 & faixa de $x$: 0 (baixa), 1 (neutro) ou 2 (alta) \\
  $\tau_1, \tau_2$          & limiares que separam as faixas (tercis do treino) \\
  $a_t$                     & ação: posição escolhida no dia $t$ para o dia seguinte, $-1$, $0$ ou $+1$ \\
  $h_{t-1}$                 & posição anterior (a ação do dia $t-1$) \\
  $s_t$                     & estado: $(\phi(x_t), \phi(x_{t-1}), h_{t-1})$ \\
  $\kappa$                  & corretagem por unidade de posição trocada \\
  $R_{t+1}$                 & recompensa: log-retorno da carteira, líquido da corretagem \\
  $W_t$                     & patrimônio no dia $t$ \\
  $\gamma$                  & fator de desconto \\
  $G_t$                     & retorno descontado a partir de $t$ \\
  $V(s)$, $Q(s,a)$          & valor do estado; valor de tomar a ação $a$ no estado $s$ \\
  $V^*$, $Q^*$, $\pi^*$     & valores ótimos e política ótima \\
  $\widehat{P}$, $\widehat{R}$ & probabilidades de transição e recompensas estimadas no treino \\
  $N(s,a)$                  & número de visitas ao par $(s,a)$ \\
  $\varepsilon$             & probabilidade de explorar (ação sorteada) \\
  $H$                       & número de dias atualizados em cada episódio \\
  \bottomrule
\end{tabular}
\end{center}

\newpage

% =============================================================
\section{O problema}

Todo dia, no fechamento (21h de Brasília, 0h UTC), o robô escolhe a posição em BTC-USD para o dia seguinte: comprado ($+1$), fora do mercado ($0$) ou vendido ($-1$). O Bitcoin negocia 24 horas por dia, 7 dias por semana; por isso, nas métricas, 1 ano $=$ 365 dias.

\textit{Trading} não é só prever se o preço sobe ou cai. Toda troca de posição paga uma corretagem de $\kappa = 0{,}1\%$ \textbf{por unidade de posição trocada}:

\begin{center}
\begin{tabular}{l c c}
  \toprule
  \textbf{Movimento} & $|a_t - h_{t-1}|$ & \textbf{Custo} \\
  \midrule
  Mantém a posição                         & 0 & 0\% \\
  Entra ou sai ($0 \leftrightarrow \pm 1$)  & 1 & 0,1\% \\
  Inverte ($+1 \leftrightarrow -1$)         & 2 & 0,2\% \\
  \bottomrule
\end{tabular}
\end{center}

Por isso a melhor ação depende da posição anterior, $h_{t-1}$, que entra no estado.

\bloco{Dados}
\begin{itemize}
  \item Preços diários do Yahoo Finance (biblioteca \texttt{yfinance}).
  \item \textbf{Treino:} 06/11/2014 a 31/12/2021. \textbf{Teste:} 01/01/2022 a 18/11/2026 (1.783 dias).
  \item Tudo é estimado só no treino. O teste é usado uma única vez, no fim.
  \item A data final é fixa no notebook (\texttt{FIM}): rodar de novo dá os mesmos números.
\end{itemize}

% =============================================================
\section{A modelagem: estado, ação e recompensa}

\bloco{Estado}
O preço do Bitcoin foi de algumas centenas de dólares em 2014 para mais de US\$ 120 mil em 2025; usar o preço bruto impediria o robô de comparar épocas. Por isso o estado usa um sinal relativo:
\[
  x_t = \ln\frac{c_t}{\mathrm{MA}_{50}(c_t)}, \qquad s_t = (x_t,\, x_{t-1},\, h_{t-1}).
\]
\begin{itemize}
  \item $x_t$ mede quanto o preço está acima ou abaixo da média móvel de 50 dias. É comparável em qualquer época.
  \item $x_{t-1}$, o valor de ontem, mostra a direção do movimento (por exemplo, cruzando a média para cima).
  \item $h_{t-1}$, a posição anterior, é necessária para calcular a corretagem.
\end{itemize}
Como $x_t$ é contínuo, o estado ainda precisa ser discretizado (Seção~\ref{sec:disc}).

\bloco{Ação}
$a_t \in \mathcal{A} = \{-1, 0, +1\}$: a posição para o dia seguinte.

\bloco{Recompensa}
\[
  R_{t+1} = \ln\!\big(1 + a_t\, r_{t+1} - \kappa\, |a_t - h_{t-1}|\big), \qquad r_{t+1} = \frac{c_{t+1}}{c_t} - 1.
\]
Duas escolhas importantes:
\begin{itemize}
  \item \textbf{Log.} O patrimônio evolui como $W_{t+1} = W_t\,(1 + a_t r_{t+1} - \kappa|a_t - h_{t-1}|)$, então $R_{t+1} = \ln(W_{t+1}/W_t)$, e a soma das recompensas é telescópica:
  \[
    \sum_{t=0}^{T-1} R_{t+1} = \ln\frac{W_1}{W_0} + \ln\frac{W_2}{W_1} + \dots + \ln\frac{W_T}{W_{T-1}} = \ln\frac{W_T}{W_0}.
  \]
  Os retornos \textbf{multiplicam} o patrimônio, e o aprendizado por reforço \textbf{soma} recompensas; o log transforma um no outro de forma exata. Maximizar $\mathbb{E}[\ln W_T]$ é o critério de Kelly (crescimento composto do capital).
  \item \textbf{Retorno simples dentro do log.} Para o vendido, o ganho é $\ln(1 - r)$, e não $-\ln(1 + r)$. Se o BTC cai 10\%, o vendido ganha $\ln(1{,}10) = 9{,}53\%$ em log; usar $-\ln(0{,}90) = 10{,}54\%$ superestimaria o ganho em 1 ponto.
\end{itemize}

\begin{exemplo}{aditividade do log em 2 dias (números hipotéticos)}
Patrimônio inicial de US\$ 1.000,00, $\kappa = 0{,}001$, começando fora do mercado ($h_0 = 0$).

\textbf{Dia 1.} O robô compra ($a_1 = +1$) e paga $\kappa\,|1 - 0| = 0{,}001$. O BTC sobe 4\%:
\[
  R_2 = \ln(1 + 0{,}04 - 0{,}001) = \ln(1{,}039) \approx 0{,}038259, \qquad W_2 = \text{US\$ } 1.039{,}00.
\]
\textbf{Dia 2.} O robô inverte para vendido ($a_2 = -1$) e paga $\kappa\,|-1 - 1| = 0{,}002$. O BTC cai 5\%:
\[
  R_3 = \ln(1 + 0{,}05 - 0{,}002) = \ln(1{,}048) \approx 0{,}046884, \qquad W_3 = 1.039 \times 1{,}048 = \text{US\$ } 1.088{,}87.
\]
\textbf{Soma:} $R_2 + R_3 = 0{,}085143 = \ln(1.088{,}87 / 1.000)$. A aditividade é exata.
\end{exemplo}

\bloco{O MDP do projeto}
Um processo de decisão de Markov (MDP) é a tupla $(\mathcal{S}, \mathcal{A}, P, R, \gamma)$:
\begin{itemize}
  \item $\mathcal{S}$: os 27 estados discretos; $\mathcal{A} = \{-1, 0, +1\}$.
  \item $P(s' \mid s,a)$: probabilidades de transição, estimadas nos dados (Seção~\ref{sec:vi}).
  \item $R$: o log-retorno líquido; $\gamma = 0{,}8$.
  \item O agente é o robô; o ambiente, o mercado. A política é uma tabela com uma ação para cada um dos 27 estados.
\end{itemize}

\bloco{Limitações da modelagem}
\begin{itemize}
  \item O estado não tem volatilidade nem volume.
  \item O \textit{short} é simétrico ao \textit{long}: não há margem, liquidação, \textit{funding} nem aluguel.
  \item A execução é no fechamento, sem \textit{slippage}.
\end{itemize}

% =============================================================
\section{Fundamentos: Markov, desconto e Bellman}

\subsection{Propriedade de Markov}

O próximo estado depende só do estado atual, e não de como se chegou nele:
\[
  P(S_{t+1} = s' \mid S_t, S_{t-1}, \dots, S_0) = P(S_{t+1} = s' \mid S_t).
\]
Por isso o estado inclui a faixa de ontem e a posição anterior: o ideal é que ele concentre toda a informação relevante. Com preços reais, isso vale só de forma aproximada; nada garante que 27 estados resumam todo o passado do mercado.

\subsection{Retorno e fator de desconto}

O retorno a partir de $t$ é a soma descontada das recompensas futuras. Pondo $\gamma$ em evidência a partir do segundo termo, ele fica recursivo:
\[
  G_t = R_{t+1} + \gamma R_{t+2} + \gamma^2 R_{t+3} + \dots = \sum_{k=0}^{\infty} \gamma^k R_{t+k+1} = R_{t+1} + \gamma\, G_{t+1}.
\]

\bloco{O que $\gamma = 0{,}8$ significa}
\begin{itemize}
  \item A recompensa de daqui a 5 dias pesa $0{,}8^5 = 32{,}8\%$. Com $\gamma = 0{,}3$, pesaria só 0,2\%.
  \item A soma dos pesos é a série geométrica: com $S = \sum_{k \ge 0} \gamma^k$, vale $S - \gamma S = 1$, logo $S = 1/(1 - \gamma) = 5$. É o ``horizonte efetivo'' de 5 dias; com $\gamma = 0{,}3$, seria 1,4 dia.
  \item Se $|R| \le R_{\max}$, então $|G_t| \le R_{\max}/(1 - \gamma)$: com $\gamma < 1$, o retorno é sempre finito.
\end{itemize}

\bloco{Por que descontar}
Conveniência matemática (o retorno fica finito e o \textit{Value Iteration} converge), incerteza sobre o futuro e, com recompensa financeira, o fato de que a recompensa imediata pode valer mais. Com $\gamma = 1$, todas as recompensas pesam igual; com $\gamma$ perto de 0, só a imediata importa.

\bloco{De onde vêm $H = 42$ e o episódio de 84 dias}
O peso de tudo o que vem depois de $H$ passos é $\sum_{k \ge H} \gamma^k = \gamma^H/(1 - \gamma)$, uma fração $\gamma^H$ do total. O notebook escolhe o menor $H$ com $\gamma^H < 10^{-4}$:
\[
  H = \left\lceil \frac{\ln 10^{-4}}{\ln \gamma} \right\rceil = \left\lceil \frac{-9{,}210}{-0{,}223} \right\rceil = \lceil 41{,}3 \rceil = 42,
  \qquad 0{,}8^{42} \approx 8{,}5 \times 10^{-5}.
\]
O episódio tem $2H = 84$ dias: os $H$ primeiros são atualizados, e os outros $H$ só completam os retornos. A mesma regra dá os episódios de 16, 28, 52, 84, 176 e 360 dias da bateria de $\gamma$ (Seção~\ref{sec:gamma}).

\subsection{Funções de valor e a equação de Bellman}

As funções de valor são esperanças do retorno:
\[
  V^\pi(s) = \mathbb{E}_\pi[G_t \mid S_t = s], \qquad Q^\pi(s,a) = \mathbb{E}_\pi[G_t \mid S_t = s, A_t = a].
\]

\bloco{Dedução da equação de Bellman}
Usando $G_t = R_{t+1} + \gamma G_{t+1}$ e condicionando na ação $a$ e no próximo estado $s'$:
\[
  \begin{aligned}
    V^\pi(s) &= \mathbb{E}_\pi[R_{t+1} + \gamma\, G_{t+1} \mid S_t = s] \\
             &= \sum_a \pi(a \mid s) \sum_{s'} P(s' \mid s,a)\,\Big[R(s,a,s') + \gamma\, \mathbb{E}_\pi[G_{t+1} \mid s, a, s']\Big] \\
             &= \sum_a \pi(a \mid s) \sum_{s'} P(s' \mid s,a)\,\big[R(s,a,s') + \gamma\, V^\pi(s')\big].
  \end{aligned}
\]
Na última linha, $\mathbb{E}_\pi[G_{t+1} \mid S_t = s, A_t = a, S_{t+1} = s']$ virou $V^\pi(s')$, que só depende de $s'$. Isso vale se o futuro depende do passado só pelo estado atual: \textbf{é aqui que entra a propriedade de Markov}.

Em termos de $Q$:
\[
  Q^\pi(s,a) = \sum_{s'} P(s' \mid s,a)\,\big[R(s,a,s') + \gamma\, V^\pi(s')\big], \qquad
  V^\pi(s) = \sum_a \pi(a \mid s)\, Q^\pi(s,a).
\]
Em palavras: o valor de uma ação é a recompensa de hoje mais o valor de onde o robô acorda amanhã, descontado por $\gamma$.

\begin{exemplo}{valor de comprar em alta (números hipotéticos, $\gamma = 0{,}8$)}
\begin{itemize}
  \item Com 70\% de chance continua em alta ($V = 0{,}10$), ganho imediato $+0{,}02$: $0{,}02 + 0{,}8 \times 0{,}10 = 0{,}100$.
  \item Com 30\% de chance vira baixa ($V = 0{,}01$), perda imediata $-0{,}03$: $-0{,}03 + 0{,}8 \times 0{,}01 = -0{,}022$.
  \item $Q(\text{alta}, +1) = 0{,}70 \times 0{,}100 + 0{,}30 \times (-0{,}022) = 0{,}0634$.
\end{itemize}
\end{exemplo}

\bloco{Forma matricial}
Para uma política fixa, empilhando os 27 estados num vetor, a equação é linear:
\[
  V^\pi = R^\pi + \gamma P^\pi V^\pi \;\Longrightarrow\; (I - \gamma P^\pi)\,V^\pi = R^\pi \;\Longrightarrow\; V^\pi = (I - \gamma P^\pi)^{-1} R^\pi.
\]
A inversa existe porque $\gamma < 1$: $(I - \gamma P^\pi)^{-1} = \sum_{k \ge 0} \gamma^k (P^\pi)^k$, ou seja,
\[
  V^\pi = R^\pi + \gamma\, P^\pi R^\pi + \gamma^2 (P^\pi)^2 R^\pi + \dots,
\]
a recompensa esperada de hoje, mais a de amanhã descontada, e assim por diante. O custo de inverter é $O(n^3)$; com 27 estados, é instantâneo.

\bloco{Otimalidade}
O valor ótimo é o melhor entre todas as políticas, $V^*(s) = \max_\pi V^\pi(s)$, e satisfaz:
\[
  V^*(s) = \max_a Q^*(s,a), \qquad
  Q^*(s,a) = \sum_{s'} P(s' \mid s,a)\,\Big[R(s,a,s') + \gamma \max_{a'} Q^*(s',a')\Big].
\]
Existe sempre uma política ótima determinística, $\pi^*(s) = \arg\max_a Q^*(s,a)$: basta uma tabela com uma ação por estado. Por causa do $\max$, essa equação não é linear e não dá para resolvê-la invertendo uma matriz. Daí o \textit{Value Iteration}.

\begin{defesa}
A dedução mostra onde está a hipótese de Markov. O \textit{Value Iteration} e o Q-learning usam a equação de Bellman, então dependem dela. O Monte Carlo estima $\mathbb{E}[G_t \mid S_t = s, A_t = a]$ direto, pela média dos retornos observados, sem esse passo. Por isso o MC \textbf{depende menos} de Markov; mas a política dele também só enxerga os 27 estados.
\end{defesa}

% =============================================================
\section{Discretização: tercis e 27 estados}
\label{sec:disc}

Como $x_t$ é contínuo, há infinitos estados, e $V(s)$ e $Q(s,a)$ não cabem numa tabela. A solução: os 2.614 valores de $x_t$ do treino são ordenados e cortados em 3 faixas com 1/3 dos dias cada. Os cortes ficam congelados e são aplicados ao teste sem recalcular, para não vazar informação do futuro.

\bloco{Limiares}
\begin{itemize}
  \item $\tau_1 = -0{,}0234$ (percentil 33\%): preço cerca de 2,3\% abaixo da média.
  \item $\tau_2 = 0{,}0953$ (percentil 67\%): preço cerca de 10\% acima da média.
  \item Não são simétricos porque, de 2014 a 2021, o Bitcoin passou mais tempo acima da média de 50 dias.
\end{itemize}

\figura{limiares_tercis.png}{0.62\textwidth}{Distribuição de $x_t$ no treino e os limiares $\tau_1$ e $\tau_2$.}

\bloco{Fórmula}
\[
  \phi(x) =
  \begin{cases}
    0 \ (\text{baixa}),  & x \le \tau_1 \\
    1 \ (\text{neutro}), & \tau_1 < x \le \tau_2 \\
    2 \ (\text{alta}),   & x > \tau_2
  \end{cases}
  \qquad
  \tilde{s}_t = \big(\phi(x_t),\, \phi(x_{t-1}),\, h_{t-1}\big).
\]
São $3 \times 3 \times 3 = 27$ estados e $27 \times 3 = 81$ pares estado-ação. Na transição, a parte de mercado vem dos dados e a posição vem da ação: $\tilde{s}_{t+1} = (\phi(x_{t+1}), \phi(x_t), a_t)$. Os três métodos usam a mesma discretização.

\bloco{Limitações}
A quantização trata igual uma queda leve ($x_t = -0{,}025$) e um colapso ($x_t = -0{,}60$). O padrão $(0,2)$, uma queda de alta para baixa em um só dia, nunca ocorre no treino.

% =============================================================
\section{Modelo estimado e \textit{Value Iteration}}
\label{sec:vi}

\subsection{Estimação de \texorpdfstring{$\widehat{P}$}{P} e \texorpdfstring{$\widehat{R}$}{R}}

O robô não mexe no preço do Bitcoin. Então só a parte de mercado da transição é aleatória: basta contar, no treino, quantas vezes cada par (hoje, ontem) foi seguido de cada faixa amanhã. A posição seguinte é a própria ação.

\bloco{Fórmulas}
Com $i = \phi(x_t)$, $k = \phi(x_{t-1})$ e $j = \phi(x_{t+1})$:
\[
  \hat{p}(j \mid i,k) = \frac{N(i,k,j)}{\sum_{j'} N(i,k,j')}, \qquad
  \widehat{P}(s' \mid s,a) = \begin{cases} \hat{p}(j \mid i,k) & \text{se } s = (i,k,h) \text{ e } s' = (j,i,a) \\ 0 & \text{caso contrário} \end{cases}
\]
\[
  \widehat{R}(s,a,s') = \frac{1}{N(i,k,j)} \sum_{t} \ln\!\big(1 + a\, r_{t+1} - \kappa\, |a - h|\big).
\]

\bloco{Por que contar é o estimador certo}
\begin{itemize}
  \item \textbf{Probabilidade condicional:} $P(A \mid B) = P(A \cap B)/P(B)$ restringe o espaço amostral aos casos em que $B$ aconteceu. É o que $\hat{p}(j \mid i,k)$ faz: entre os dias com o par $(i,k)$, conta a fração seguida da faixa $j$.
  \item \textbf{Máxima verossimilhança:} para cada par $(i,k)$, a faixa de amanhã é uma variável categórica, e o estimador de máxima verossimilhança é exatamente a frequência relativa.
  \item \textbf{Lei dos Grandes Números:} a frequência relativa é uma média de indicadoras ($\mathbb{E}[\mathbf{1}_A] = P(A)$) e converge para a probabilidade real. O erro cai como $\sigma/\sqrt{n}$, então $\widehat{P}$ e $\widehat{R}$ são confiáveis nas transições frequentes e frágeis nas raras, como $(2,0)$. Os dias não são independentes, mas a convergência continua valendo numa cadeia ergódica.
\end{itemize}

\begin{center}
\small
\begin{tabular}{c c ccc}
  \toprule
  $(i,k) = (\phi(x_t), \phi(x_{t-1}))$ & Regime & $j = 0$ & $j = 1$ & $j = 2$ \\
  \midrule
  (0,0) & Baixa estável        & 0,933 & 0,065 & 0,001 \\
  (0,1) & Neutro $\to$ baixa   & 0,915 & 0,085 & 0,000 \\
  (0,2) & Alta $\to$ baixa     & 0,333 & 0,333 & 0,333 \\
  \midrule
  (1,0) & Baixa $\to$ neutro   & 0,310 & 0,690 & 0,000 \\
  (1,1) & Neutro estável       & 0,053 & 0,889 & 0,057 \\
  (1,2) & Alta $\to$ neutro    & 0,016 & 0,672 & 0,312 \\
  \midrule
  (2,0) & Baixa $\to$ alta     & 0,000 & 1,000 & 0,000 \\
  (2,1) & Neutro $\to$ alta    & 0,000 & 0,333 & 0,667 \\
  (2,2) & Alta estável         & 0,000 & 0,052 & 0,948 \\
  \bottomrule
\end{tabular}
\end{center}

\bloco{Leitura}
As faixas são persistentes: a baixa estável continua em baixa em 93,3\% dos dias, e a alta estável, em 94,8\%. O padrão $(0,2)$ nunca ocorre no treino; nele se usa a distribuição uniforme e retorno nulo.

\begin{exemplo}{quanto dura um regime (supondo Markov)}
Na baixa estável $(0,0)$, amanhã continua em baixa com $p = 0{,}933$, e o par volta a ser $(0,0)$.
\begin{itemize}
  \item Continuar em baixa por mais 2 dias: $0{,}933^2 \approx 0{,}870$. Por mais 5 dias: $0{,}933^5 \approx 0{,}707$.
  \item O número de dias seguidos no regime segue uma distribuição geométrica, com média $1/(1 - p) = 1/0{,}067 \approx 15$ dias. Na alta estável, $1/(1 - 0{,}948) \approx 19$ dias.
\end{itemize}
Os regimes duram semanas: por isso olhar a faixa de hoje e a de ontem ajuda a decidir.
\end{exemplo}

\subsection{\textit{Value Iteration}}

Começa com $V_0(s) = 0$ em todos os estados. A cada rodada, cada estado escolhe a ação que maximiza a recompensa esperada mais $\gamma$ vezes o valor do próximo estado; o valor do futuro se propaga um passo para trás por rodada, até estabilizar:
\[
  V_{k+1}(s) = \max_{a} \sum_{s'} \widehat{P}(s' \mid s,a)\,\big[\widehat{R}(s,a,s') + \gamma\, V_k(s')\big], \qquad
  \pi^*(s) = \arg\max_{a} Q^*(s,a).
\]
O critério de parada é $\max_s |V_{k+1}(s) - V_k(s)| < \varepsilon$. No projeto, com $\varepsilon = 10^{-12}$, foram 98 iterações; o VI chegou a 99\% do valor ótimo já na iteração 15.

\begin{exemplo}{duas iterações à mão (números hipotéticos, $\gamma = 0{,}8$)}
No estado $A$, comprar leva a $A$ com 70\% de chance (recompensa $+0{,}02$) e a $B$ com 30\% (recompensa $-0{,}03$, $V(B) = 0$). Ficar fora dá 0.

\textbf{Iteração 1} (partindo de $V_0 = 0$):
\[
  Q_1(A,+1) = 0{,}70 \times 0{,}02 + 0{,}30 \times (-0{,}03) = 0{,}0050 \;\Rightarrow\; V_1(A) = 0{,}0050
\]
\textbf{Iteração 2:}
\[
  \begin{aligned}
    Q_2(A,+1) &= 0{,}70 \times (0{,}02 + 0{,}8 \times 0{,}0050) + 0{,}30 \times (-0{,}03) \\
              &= 0{,}0168 - 0{,}0090 = 0{,}0078 \;\Rightarrow\; V_2(A) = 0{,}0078
  \end{aligned}
\]
\end{exemplo}

\bloco{Por que converge}
O VI aplica repetidamente o operador $(TV)(s) = \max_a \sum_{s'} P(s' \mid s,a)\,[R(s,a,s') + \gamma V(s')]$.
\begin{enumerate}
  \item \textbf{Lema:} $\big|\max_a f(a) - \max_a g(a)\big| \le \max_a |f(a) - g(a)|$. Se $\max f \ge \max g$ e $a^*$ é o argmax de $f$, então $\max f - \max g \le f(a^*) - g(a^*) \le \max_a |f(a) - g(a)|$.
  \item \textbf{Contração:} com o lema, as recompensas se cancelam e só sobra $\gamma$:
  \[
    \begin{aligned}
      |(TV_1)(s) - (TV_2)(s)| &\le \max_a \Big|\gamma \sum_{s'} P(s' \mid s,a)\,\big(V_1(s') - V_2(s')\big)\Big| \\
                              &\le \gamma \max_a \sum_{s'} P(s' \mid s,a)\, \|V_1 - V_2\|_\infty = \gamma\, \|V_1 - V_2\|_\infty.
    \end{aligned}
  \]
  \item \textbf{Ponto fixo de Banach:} existe um único $V^*$ com $TV^* = V^*$, e $V_k \to V^*$ partindo de qualquer $V_0$. O erro cai pelo menos pelo fator $\gamma$ a cada iteração: $\|V_k - V^*\|_\infty \le \gamma^k\, \|V_0 - V^*\|_\infty$.
  \item \textbf{Garantia do critério de parada:} se $\|V_{k+1} - V_k\| < \varepsilon$, então
  \[
    \|V_{k+1} - V^*\| \le \gamma\, \|V_k - V^*\| \le \gamma\,\big(\|V_k - V_{k+1}\| + \|V_{k+1} - V^*\|\big)
    \;\Longrightarrow\; \|V_{k+1} - V^*\| < \frac{\gamma\,\varepsilon}{1 - \gamma}.
  \]
  Com $\gamma = 0{,}8$ e $\varepsilon = 10^{-12}$, o erro final é menor que $4 \times 10^{-12}$.
\end{enumerate}

\subsection{\textit{Policy Iteration}, a alternativa}

Alterna dois passos: \textbf{avaliar} $\pi$ (resolvendo o sistema linear da forma matricial) e \textbf{melhorar}, com $\pi'(s) = \arg\max_a Q^\pi(s,a)$. Para quando a política não muda mais.

\bloco{Teorema da melhoria}
Se $Q^\pi(s, \pi'(s)) \ge V^\pi(s)$ em todo $s$, então $V^{\pi'}(s) \ge V^\pi(s)$ em todo $s$. A prova aplica a hipótese um passo de cada vez:
\[
  \begin{aligned}
    V^\pi(s) &\le Q^\pi(s, \pi'(s)) = \mathbb{E}_{\pi'}\big[R_{t+1} + \gamma V^\pi(S_{t+1}) \mid S_t = s\big] \\
             &\le \mathbb{E}_{\pi'}\big[R_{t+1} + \gamma R_{t+2} + \gamma^2 V^\pi(S_{t+2}) \mid S_t = s\big]
              \le \dots \le V^{\pi'}(s).
  \end{aligned}
\]
Como cada melhoria aumenta o valor ou para, nenhuma política se repete; e há um número finito de políticas determinísticas ($3^{27} \approx 7{,}6 \times 10^{12}$). Logo o PI termina. VI e PI chegam ao mesmo $V^*$.

\subsection{Política ótima do VI}

\begin{center}
\small
\begin{tabular}{c l ccc}
  \toprule
  $(i,k)$ & Regime & $h_{t-1} = -1$ & $h_{t-1} = 0$ & $h_{t-1} = +1$ \\
  \midrule
  (0,0) & Baixa estável      & $0$  & $0$  & $+1$ \\
  (0,1) & Neutro $\to$ baixa & $-1$ & $-1$ & $-1$ \\
  (0,2) & Alta $\to$ baixa   & $-1$ & $0$  & $+1$ \\
  \midrule
  (1,0) & Baixa $\to$ neutro & $-1$ & $0$  & $0$  \\
  (1,1) & Neutro estável     & $0$  & $0$  & $+1$ \\
  (1,2) & Alta $\to$ neutro  & $+1$ & $+1$ & $+1$ \\
  \midrule
  (2,0) & Baixa $\to$ alta   & $-1$ & $-1$ & $-1$ \\
  (2,1) & Neutro $\to$ alta  & $+1$ & $+1$ & $+1$ \\
  (2,2) & Alta estável       & $+1$ & $+1$ & $+1$ \\
  \bottomrule
\end{tabular}
\end{center}

\begin{itemize}
  \item \textbf{Baixa estável:} fica fora, ou segue no \textit{long} se já estava comprado.
  \item \textbf{Entrando na baixa} $(0,1)$: vende.
  \item \textbf{Neutro estável:} só mantém o \textit{long} já aberto; sem posição, fica fora.
  \item \textbf{Alta e alta $\to$ neutro} $(2,2)$, $(2,1)$, $(1,2)$: compra.
  \item \textbf{Baixa $\to$ alta} $(2,0)$: vende, mas esse salto é raro no treino.
  \item \textbf{Padrão $(0,2)$}, que não ocorre no treino: mantém a posição, porque sem dados a única recompensa é o custo da troca.
\end{itemize}
Nos estados estáveis, ``manter o \textit{long} já aberto'' reflete a corretagem (trocar custa) e o viés de alta do treino.

% =============================================================
\section{Métricas de avaliação}
\label{sec:metricas}

Comparar estratégias só pelo retorno total engana: uma estratégia pode dobrar o capital passando por uma queda de 80\% no caminho. Por isso usamos várias métricas, e cada uma olha um lado do resultado.

\begin{center}
\small
\begin{tabular}{L{2.6cm} L{9.2cm} L{2.9cm}}
  \toprule
  \textbf{Métrica} & \textbf{O que responde} & \textbf{Melhor quando} \\
  \midrule
  Retorno total & Quanto o dinheiro cresceu do início ao fim? & maior \\
  Retorno anual & Quanto isso dá por ano? & maior \\
  Sharpe        & Quanto se ganhou acima da renda fixa para cada unidade de oscilação? & maior \\
  MDD           & Qual foi a pior queda do patrimônio, de um pico até o fundo? & mais perto de 0 \\
  Calmar        & Quanto se ganhou por ano para cada ponto dessa pior queda? & maior \\
  Trocas        & Quantas vezes a posição mudou e pagou corretagem? & menor \\
  \bottomrule
\end{tabular}
\end{center}

\bloco{Retorno total e retorno anual}
O retorno total é quanto o patrimônio cresceu, $W_T/W_0 - 1$. No teste, o VI fez 127,80\%: US\$ 1.000 viraram cerca de US\$ 2.278.

O retorno anual converte isso numa taxa por ano, como se o patrimônio tivesse crescido igual todo ano. Serve para comparar períodos de tamanhos diferentes: o treino tem 7 anos, e o teste, quase 5.
\[
  (1 + R)^{T/365} = \frac{W_T}{W_0} \;\Longrightarrow\; R_{\text{anual}} = \Big(\frac{W_T}{W_0}\Big)^{365/T} - 1.
\]
Os 127,80\% do VI em quase 5 anos equivalem a 18,37\% ao ano.

\bloco{Índice de Sharpe}
Mede o retorno \textbf{por unidade de risco}. O numerador é quanto a estratégia rendeu acima do que se ganharia sem risco; o denominador, quanto o patrimônio oscila de um dia para o outro. Entre duas estratégias com o mesmo retorno, a que oscila menos tem Sharpe maior.
\[
  S = \sqrt{365}\,\frac{\bar{e}}{\sigma_e}, \qquad e_t = \ln\frac{W_t}{W_{t-1}} - r^f_t, \qquad r^f_t = \frac{\ln(1 + i_t)}{365}.
\]
\begin{itemize}
  \item $e_t$ é o retorno do dia acima da taxa livre de risco; $\bar e$ é a média desses retornos, e $\sigma_e$, o desvio padrão.
  \item A taxa livre $i_t$ é o T-bill americano de 3 meses, porque o BTC é cotado em dólar: média de 0,78\% a.a. no treino e 3,94\% a.a. no teste.
  \item \textbf{Como ler:} Sharpe 0 quer dizer que a estratégia rendeu o mesmo que a renda fixa; negativo, que rendeu menos; quanto maior, melhor. No teste, o \textit{Buy \& Hold} subiu 31,95\%, mas oscilou tanto que o Sharpe ficou em 0,04; o do VI foi 0,46.
  \item \textbf{Por que $\sqrt{365}$:} se os $e_t$ forem i.i.d., em 365 dias a média se soma ($365\,\bar e$) e a variância também, então o desvio cresce com a raiz ($\sqrt{365}\,\sigma_e$). A razão anual fica $365\,\bar e/(\sqrt{365}\,\sigma_e) = \sqrt{365}\,\bar e/\sigma_e$. Em ações, que negociam cerca de 252 dias por ano, seria $\sqrt{252}$.
\end{itemize}

\bloco{\textit{Maximum drawdown} (MDD)}
A pior queda do patrimônio, medida de um pico até o fundo seguinte. Responde: ``quanto eu teria perdido se tivesse entrado no pior momento?''
\[
  \mathrm{MDD} = \min_t \left(\frac{W_t}{P_t} - 1\right), \qquad P_t = \max_{u \le t} W_u \ \ \text{(o maior patrimônio até o dia $t$)}.
\]
No teste, o MDD do \textit{Buy \& Hold} foi $-66{,}93\%$: em algum momento, o patrimônio caiu para cerca de um terço do pico anterior. O do VI foi $-29{,}35\%$.

\bloco{Índice de Calmar}
O retorno anual dividido pelo tamanho da pior queda: quanto se ganha por ano para cada ponto de queda que foi preciso aguentar.
\[
  \text{Calmar} = \frac{R_{\text{anual}}}{|\mathrm{MDD}|}.
\]
No teste, o VI tem $18{,}37/29{,}35 \approx 0{,}63$; o \textit{Buy \& Hold}, $5{,}84/66{,}93 \approx 0{,}09$. O Sharpe pune toda oscilação; o Calmar pune só o pior tombo.

\bloco{Trocas}
Quantas vezes a posição mudou. Cada troca paga corretagem, então, para o mesmo resultado, menos trocas é melhor. No teste, o Q-learning trocou 190 vezes, e o VI, 113.

\bloco{Por que essas métricas fazem sentido aqui}
\begin{itemize}
  \item O retorno anual é o próprio objetivo do agente: a soma das recompensas é o log do patrimônio final.
  \item O Sharpe compara estratégias com exposições muito diferentes (o RL fica fora boa parte do tempo, o \textit{Buy \& Hold} nunca sai).
  \item O MDD é essencial no Bitcoin: o \textit{Buy \& Hold} caiu 83\% no treino e 67\% no teste.
  \item O número de trocas importa porque a corretagem está na recompensa.
  \item Os modelos econométricos não fazem o mesmo tipo de previsão que o RL; comparar todos pelo patrimônio, com a mesma corretagem, é a única base comum.
  \item No simulador, o caixa não rende juros, mas o Sharpe desconta a taxa livre. Isso penaliza quem fica fora, ou seja, o RL. A escolha é conservadora.
\end{itemize}

\bloco{Erro amostral do Sharpe}
O Sharpe é uma estimativa. Com retornos i.i.d. e $T$ dias:
\[
  \mathrm{ep}(\hat S_{\text{diário}}) \approx \sqrt{\frac{1 + \hat S_{\text{diário}}^2/2}{T}}
  \quad\Longrightarrow\quad
  \mathrm{ep}(\hat S_{\text{anual}}) \approx \sqrt{\frac{365}{T}}.
\]
No teste, $T = 1.783$ dias: o erro padrão do Sharpe anual é cerca de $\sqrt{365/1.783} \approx 0{,}45$, e um intervalo aproximado de 95\% para o Sharpe do VI seria $0{,}46 \pm 0{,}89$.

\begin{defesa}
Com 5 anos de teste, \textbf{nenhuma diferença isolada de Sharpe é estatisticamente significativa}: 0,46 contra 0,44 está dentro do ruído. O argumento do trabalho é a consistência: os três RL acima de todos os econométricos em Sharpe e Calmar, com queda máxima menor que a de todos eles, e resultados parecidos com 5 sementes. Comparar dois Sharpe formalmente exigiria considerar a correlação entre as estratégias (teste de Jobson--Korkie), o que não fizemos.
\end{defesa}

% =============================================================
\section{Monte Carlo}

Estimar $\widehat{P}$ supõe que o mercado segue um modelo de Markov discreto e estável. O Monte Carlo aprende $Q(s,a)$ direto do histórico, sem montar o modelo (\textit{model-free}).

\bloco{Como funciona}
\begin{itemize}
  \item 30.000 episódios de 84 dias sobre o histórico real do treino.
  \item Cada episódio começa num dia sorteado, com posição inicial sorteada e primeira ação aleatória (\textit{exploring starts}).
  \item No fim do episódio, calcula o retorno $G_t$ de cada passo, de trás para frente, e move a média $Q(s,a)$ na direção dele.
  \item Só os primeiros $H = 42$ dias são atualizados, para que todo $G_t$ tenha pelo menos 42 recompensas futuras.
\end{itemize}

\bloco{Fórmulas}
\[
  G_t = R_{t+1} + \gamma\, G_{t+1}, \qquad
  Q(s,a) \leftarrow Q(s,a) + \frac{1}{N(s,a)}\big[G_t - Q(s,a)\big], \qquad
  \pi(s) = \arg\max_a Q(s,a).
\]
O passo $1/N$ faz de $Q(s,a)$ a média simples dos retornos. É a média incremental:
\[
  M_n = \frac{1}{n}\sum_{i=1}^{n} X_i = \frac{1}{n}\Big(X_n + (n-1)\,M_{n-1}\Big) = M_{n-1} + \frac{1}{n}\big(X_n - M_{n-1}\big).
\]

\begin{exemplo}{uma atualização (números hipotéticos)}
$Q(s,a) = 0{,}010$ depois de 4 visitas. Num novo episódio, as recompensas seguintes são $+0{,}02$ e $+0{,}03$:
\begin{itemize}
  \item Retorno: $G_t = 0{,}02 + 0{,}8 \times 0{,}03 = 0{,}044$.
  \item Quinta visita: $N = 5$, passo $1/5 = 0{,}20$.
  \item $Q(s,a) \leftarrow 0{,}010 + 0{,}20 \times (0{,}044 - 0{,}010) = 0{,}0168$.
\end{itemize}
\end{exemplo}

\bloco{\textit{First-visit} $\times$ \textit{every-visit}}
\begin{itemize}
  \item \textbf{First-visit:} em cada episódio, usa só a primeira passagem pelo par $(s,a)$. Os retornos vêm de episódios diferentes, são independentes, e a média é sem viés. Indicado quando há forte correlação temporal.
  \item \textbf{Every-visit:} usa todas as passagens. Aproveita mais amostras e converge mais rápido, mas os retornos do mesmo episódio são correlacionados.
\end{itemize}

\begin{exemplo}{\textit{first-visit} $\times$ \textit{every-visit}}
Dois episódios com $\gamma = 1$ e recompensa só no fim: $A \to B \to A \to C \to D$, com recompensa final 5, e $A \to A \to A \to B \to D$, com recompensa final 3. Toda visita a $A$ tem retorno igual à recompensa final do seu episódio.
\begin{itemize}
  \item \textit{First-visit}: uma visita por episódio, $V(A) = (5 + 3)/2 = 4{,}0$.
  \item \textit{Every-visit}: duas visitas no 1º episódio e três no 2º, $V(A) = (5 + 5 + 3 + 3 + 3)/5 = 3{,}8$.
\end{itemize}
\end{exemplo}

\bloco{Exploração $\varepsilon$-greedy}
Com probabilidade $\varepsilon$, o robô sorteia uma ação; senão, escolhe a melhor:
\[
  \pi(a \mid s) =
  \begin{cases}
    1 - \varepsilon + \varepsilon/|\mathcal{A}| & \text{para a melhor ação} \\
    \varepsilon/|\mathcal{A}|                   & \text{para as outras}
  \end{cases}
\]
Com 3 ações e $\varepsilon = 0{,}1$: 0,933 e 0,033. No projeto, $\varepsilon$ cai linearmente de 1 para 0,05 nos primeiros 60\% dos episódios. Explorar demais deixa o aprendizado lento; explorar de menos arrisca convergir cedo para uma política ruim.

\bloco{No projeto}
\begin{itemize}
  \item Com a semente 42, o \textit{every-visit} fez 1.260.000 atualizações (contra 388.071) e teve Sharpe no treino de 1,20 (contra 1,12). Por isso foi escolhido.
  \item Mesmo com \textit{exploring starts}, só 72 dos 81 pares são visitados, porque o padrão $(0,2)$ nunca ocorre.
\end{itemize}

\bloco{O que o MC estima}
\begin{itemize}
  \item \textbf{Sem viés, mas com variância alta.} $G_t$ é uma amostra exata do retorno, mas soma muitas recompensas aleatórias. Se fossem independentes, com variância $\sigma^2$ cada:
  \[
    \mathrm{Var}(G_t) = \sum_{k \ge 0} \gamma^{2k}\sigma^2 = \frac{\sigma^2}{1 - \gamma^2} = \frac{\sigma^2}{0{,}36} \approx 2{,}8\,\sigma^2 \qquad (\gamma = 0{,}8).
  \]
  \item \textbf{Mistura de políticas.} Com passo $1/N$, $Q$ é a média de todos os retornos, inclusive os do começo do treino, quando $\varepsilon$ estava perto de 1. Então o $Q$ do MC não é exatamente $Q^*$; para o controle MC chegar à política ótima, $\varepsilon$ teria de ir a zero, e o nosso para em 0,05. Isso ajuda a explicar por que o MC fica mais longe de $Q^*$ que o Q-learning, mesmo com uma política melhor no teste.
\end{itemize}

% =============================================================
\section{Q-learning}

O Monte Carlo só atualiza no fim do episódio. O Q-learning atualiza a cada dia, usando a estimativa do próximo estado (\textit{bootstrap}).

\bloco{Fórmulas}
O erro de diferença temporal (TD) compara a estimativa atual com ``recompensa de hoje + valor descontado de amanhã'':
\[
  \delta_t = R_{t+1} + \gamma \max_{a'} Q(s',a') - Q(s,a), \qquad
  Q(s,a) \leftarrow Q(s,a) + \frac{1}{N(s,a)}\,\delta_t.
\]
O alvo usa a melhor ação em $s'$, não a ação realmente tomada: o Q-learning aprende o valor da política gulosa enquanto segue a exploratória. Por isso é \textit{off-policy}. Ao usar o próximo estado, supõe, como o VI, que o estado é Markov.

\begin{exemplo}{um passo (números hipotéticos)}
$Q(s,+1) = 0{,}010$ com passo $1/5 = 0{,}20$. O robô compra, recebe $R_{t+1} = 0{,}015$ e vai para $s'$, onde $Q(s',-1) = -0{,}01$, $Q(s',0) = 0{,}005$ e $Q(s',+1) = 0{,}025$.
\begin{itemize}
  \item Melhor valor em $s'$: $\max_{a'} Q(s',a') = 0{,}025$.
  \item Alvo: $0{,}015 + 0{,}8 \times 0{,}025 = 0{,}035$.
  \item $Q(s,+1) \leftarrow 0{,}010 + 0{,}20 \times (0{,}035 - 0{,}010) = 0{,}015$.
\end{itemize}
\end{exemplo}

\bloco{SARSA $\times$ Q-learning}
O SARSA usa no alvo a ação realmente tomada, $Q(s', a')$: é \textit{on-policy} e aprende o valor da política exploratória. O Q-learning usa o máximo.

\begin{exemplo}{SARSA $\times$ Q-learning no mesmo passo}
Mesmos números. Suponha que a política $\varepsilon$-greedy, explorando, escolha em $s'$ a ação $0$, com $Q(s',0) = 0{,}005$.
\begin{itemize}
  \item SARSA: alvo $= 0{,}015 + 0{,}8 \times 0{,}005 = 0{,}019$, e $Q(s,+1) \leftarrow 0{,}010 + 0{,}20 \times 0{,}009 = 0{,}0118$.
  \item Q-learning: alvo $= 0{,}015 + 0{,}8 \times 0{,}025 = 0{,}035$, e $Q(s,+1) \leftarrow 0{,}015$.
\end{itemize}
O Q-learning aprende o valor da política gulosa, que é a que executamos no teste.
\end{exemplo}

\bloco{Monte Carlo $\times$ TD}
O MC tem alta variância e baixo viés; o TD, menor variância e maior viés (enquanto a estimativa do próximo estado está errada). A ligação entre os dois: com $V$ fixo durante o episódio e $V = 0$ no fim, somando e subtraindo $\gamma V(S_{t+1})$,
\[
  G_t - V(S_t) = \delta_t + \gamma\big(G_{t+1} - V(S_{t+1})\big) = \delta_t + \gamma\,\delta_{t+1} + \gamma^2\,\delta_{t+2} + \dots
\]
O MC corrige com a soma inteira no fim; o TD corrige com o primeiro termo a cada passo.

\bloco{Por que o Q-learning aprende $Q^*$}
Se $Q = Q^*$, o valor esperado do alvo é o próprio $Q^*$, pela equação de otimalidade:
\[
  \mathbb{E}\Big[R_{t+1} + \gamma \max_{a'} Q^*(S_{t+1}, a') \;\Big|\; s, a\Big]
  = \sum_{s'} P(s' \mid s,a)\Big[R(s,a,s') + \gamma \max_{a'} Q^*(s',a')\Big] = Q^*(s,a).
\]
$Q^*$ é o ponto em que a atualização, em média, não mexe mais. A política que gerou os dados só decide quais pares são visitados.

\bloco{Condições de convergência}
Todos os pares visitados infinitas vezes, $\gamma < 1$, estado Markov e passos com $\sum_n \alpha_n = \infty$ e $\sum_n \alpha_n^2 < \infty$. O passo $1/n$ cumpre as duas: $\sum 1/n$ (série harmônica) diverge e $\sum 1/n^2 = \pi^2/6$. Um $\alpha$ constante não cumpre a segunda: a estimativa continua oscilando, o que ajuda quando o ambiente muda, mas não converge.

\bloco{Por que ele fica perto do $Q^*$ do VI}
Com passo $1/N$, cada $Q(s,a)$ é a média dos alvos observados depois de $(s,a)$ no histórico. No limite, isso se aproxima de
\[
  Q(s,a) \approx \sum_{s'} \widehat{P}(s' \mid s,a)\Big[\widehat{R}(s,a,s') + \gamma \max_{a'} Q(s',a')\Big],
\]
a equação de otimalidade do \textbf{mesmo modelo empírico} que o VI resolve (equivalência de certeza). Se o estado fosse exatamente Markov e a amostra infinita, os dois coincidiriam.

\bloco{Superestimação}
O máximo de estimativas ruidosas tende a passar do máximo verdadeiro: $\mathbb{E}[\max_a \widehat{Q}(a)] \ge \max_a \mathbb{E}[\widehat{Q}(a)]$. Exemplo: duas ações com valor real 0 e estimativas de $+0{,}01$ ou $-0{,}01$, com 50\% de chance cada e independentes. O máximo só é $-0{,}01$ se as duas saírem negativas (25\%), então $\mathbb{E}[\max] = 0{,}75 \times 0{,}01 - 0{,}25 \times 0{,}01 = 0{,}005 > 0$.

\subsection{Os três métodos lado a lado}

\begin{center}
\small
\begin{tabular}{l ccc}
  \toprule
  & \textbf{\textit{Value Iteration}} & \textbf{Monte Carlo} & \textbf{Q-learning} \\
  \midrule
  Precisa do modelo?   & sim ($\widehat{P}$, $\widehat{R}$) & não                & não \\
  Alvo da atualização  & $\sum_{s'} \widehat{P}\,[\widehat{R} + \gamma V(s')]$ & retorno real $G_t$ & $R_{t+1} + \gamma \max_{a'} Q(s',a')$ \\
  Quando atualiza      & todos os estados                   & fim do episódio     & a cada passo \\
  Supõe Markov?        & sim                                & no alvo, não        & sim \\
  Aprende              & --                                 & \textit{on-policy}  & \textit{off-policy} \\
  \midrule
  Custo do treino      & 98 iterações                       & 1.260.000           & 2.520.000 \\
  Sharpe no treino     & 1,09                               & 1,20                & 1,37 \\
  \bottomrule
\end{tabular}

\smallskip
{\footnotesize Custo de MC e QL em atualizações de $Q$; Sharpe com a semente 42.}
\end{center}

\begin{center}
\small
\begin{tabular}{L{2.4cm} L{4.2cm} L{7.6cm}}
  \toprule
  \textbf{Método} & \textbf{Converge para} & \textbf{Condições} \\
  \midrule
  \textit{Value Iteration}  & $V^*$ e $\pi^*$ do modelo estimado & $\gamma < 1$ (contração) \\
  \textit{Policy Iteration} & $\pi^*$ do modelo estimado & estados e ações finitos \\
  Monte Carlo               & melhor política $\varepsilon$-gulosa; ótima só se $\varepsilon \to 0$ & todos os pares visitados, $\varepsilon > 0$ \\
  Q-learning                & $Q^*$ & todos os pares infinitas vezes, $\sum \alpha = \infty$, $\sum \alpha^2 < \infty$, estado Markov \\
  \bottomrule
\end{tabular}
\end{center}

% =============================================================
\section{Escolha de \texorpdfstring{$\gamma$}{gama}}
\label{sec:gamma}

$\gamma$ define o \textbf{objetivo} do agente, não é só um ajuste fino. Como a soma das recompensas é o log do patrimônio final, maximizar a carteira exige valorizar o futuro; com $\gamma$ perto de 1, o objetivo se aproxima do critério de Kelly. Com $\gamma$ baixo, o agente fica míope.

\bloco{A bateria de validação}
Tudo foi refeito para cada $\gamma$, com treino até 2018 e avaliação em 2019--2021, sem usar o teste. O horizonte $H$ acompanha $\gamma$, e Monte Carlo e Q-learning usam a média de 5 sementes.

\begin{center}
\small
\begin{tabular}{c c c cccc}
  \toprule
  $\gamma$ & Peso em 5 dias ($\gamma^5$) & Episódio (dias) & VI & MC & QL & Média \\
  \midrule
  0,3  & 0,2\%  & 16  & 1,70 & 1,67 & 1,69 & 1,69 \\
  0,5  & 3,1\%  & 28  & 1,70 & 1,67 & 1,57 & 1,65 \\
  0,7  & 16,8\% & 52  & 1,59 & 1,54 & 1,40 & 1,51 \\
  \midrule
  \textbf{0,8} & \textbf{32,8\%} & \textbf{84} & \textbf{1,59} & \textbf{1,42} & \textbf{1,20} & \textbf{1,40} \\
  0,9  & 59,0\% & 176 & 1,59 & 1,17 & 1,16 & 1,31 \\
  0,95 & 77,4\% & 360 & 1,59 & 1,18 & 1,17 & 1,31 \\
  \bottomrule
\end{tabular}
\end{center}

\begin{defesa}
\begin{itemize}
  \item O $\gamma = 0{,}8$ foi fixado antes, pelo objetivo. A bateria só confere a sensibilidade.
  \item $\gamma$ menor vai melhor na validação, mas olha só 1 a 2 dias à frente: com $\gamma = 0{,}3$, o 5º dia pesa 0,2\%.
  \item O $\gamma = 0{,}8$ já desconta forte (o 5º dia pesa 32,8\%). O que se evita é o extremo míope, que ignora que a posição de hoje define o custo de amanhã.
  \item Entre os $\gamma \ge 0{,}8$, o 0,8 tem o maior Sharpe médio. O VI quase não muda com $\gamma$.
  \item Validação $\neq$ teste: outro período e treino só até 2018, por isso o Sharpe é outro.
  \item No teste, consultado só por transparência, o Sharpe do VI é o mesmo (0,46) para todo $\gamma \le 0{,}8$, e o Q-learning é o mais sensível a $\gamma$.
\end{itemize}
\end{defesa}

% =============================================================
\section{Políticas, convergência e sementes}
\label{sec:sementes}

\bloco{Política aprendida pelos três métodos (semente 42)}
\begin{center}
\small
\setlength{\tabcolsep}{4pt}
\begin{tabular}{c l ccc ccc ccc}
  \toprule
  & & \multicolumn{3}{c}{\textbf{VI} ($h_{t-1}$)} & \multicolumn{3}{c}{\textbf{MC} ($h_{t-1}$)} & \multicolumn{3}{c}{\textbf{QL} ($h_{t-1}$)} \\
  \cmidrule(lr){3-5} \cmidrule(lr){6-8} \cmidrule(lr){9-11}
  \textbf{Estado} & \textbf{Regime} & $-1$ & $0$ & $+1$ & $-1$ & $0$ & $+1$ & $-1$ & $0$ & $+1$ \\
  \midrule
  (0,0) & Baixa estável      & $0$ & $0$ & $+1$ & $0$ & $0$ & $+1$ & $0$ & $0$ & $+1$ \\
  (0,1) & Neutro $\to$ baixa & $-1$ & $-1$ & $\dif{-1}$ & $-1$ & $-1$ & $\dif{0}$ & $-1$ & $-1$ & $\dif{0}$ \\
  (0,2) & Alta $\to$ baixa   & $\dif{-1}$ & $0$ & $\dif{+1}$ & $\dif{0}$ & $0$ & $\dif{0}$ & $\dif{0}$ & $0$ & $\dif{0}$ \\
  \midrule
  (1,0) & Baixa $\to$ neutro & $\dif{-1}$ & $\dif{0}$ & $\dif{0}$ & $\dif{0}$ & $\dif{0}$ & $\dif{-1}$ & $\dif{0}$ & $\dif{-1}$ & $\dif{0}$ \\
  (1,1) & Neutro estável     & $0$ & $\dif{0}$ & $+1$ & $0$ & $\dif{0}$ & $+1$ & $0$ & $\dif{+1}$ & $+1$ \\
  (1,2) & Alta $\to$ neutro  & $+1$ & $+1$ & $+1$ & $+1$ & $+1$ & $+1$ & $+1$ & $+1$ & $+1$ \\
  \midrule
  (2,0) & Baixa $\to$ alta   & $-1$ & $-1$ & $-1$ & $-1$ & $-1$ & $-1$ & $-1$ & $-1$ & $-1$ \\
  (2,1) & Neutro $\to$ alta  & $+1$ & $+1$ & $+1$ & $+1$ & $+1$ & $+1$ & $+1$ & $+1$ & $+1$ \\
  (2,2) & Alta estável       & $+1$ & $+1$ & $+1$ & $+1$ & $+1$ & $+1$ & $+1$ & $+1$ & $+1$ \\
  \bottomrule
\end{tabular}

\smallskip
{\footnotesize Em negrito, os estados em que os métodos discordam.}
\end{center}

\begin{itemize}
  \item Os três escolhem a mesma ação em 20 dos 27 estados; MC e QL concordam em 24.
  \item Nos regimes de alta, os três compram em $(2,1)$, $(2,2)$ e $(1,2)$ e vendem no salto $(2,0)$.
  \item As diferenças ficam perto da baixa e no neutro. No neutro estável sem posição, o Q-learning compra; VI e MC ficam fora.
\end{itemize}

\bloco{Convergência}
O $Q^*$ do VI no mesmo MDP serve como referência teórica. Na média das 5 sementes, o Q-learning chega mais perto de $Q^*$ (erro médio de 0,005, contra 0,007 do MC), mas varia mais entre as sementes. Nos dois, a política final coincide com a do VI em cerca de três quartos dos estados. O padrão $(0,2)$ nunca é visitado e não tem como ser aprendido.

\figura{curva_sementes.png}{\textwidth}{Convergência em 5 sementes: política igual à do VI (esquerda) e erro médio $|Q - Q^*|$ (direita).}

\bloco{Variabilidade entre sementes}
Os resultados principais usam a semente 42. Monte Carlo e Q-learning foram repetidos com as sementes 42 a 46:

\begin{center}
\small
\begin{tabular}{l c ccc}
  \toprule
  & \textbf{Treino} & \multicolumn{3}{c}{\textbf{Teste (2022 a 2026)}} \\
  \cmidrule(lr){2-2} \cmidrule(lr){3-5}
  \textbf{Método} & \textbf{Sharpe} & \textbf{Sharpe} & \textbf{MDD} & \textbf{Calmar} \\
  \midrule
  MC \textit{first-visit} & 1,42 (0,17) & 0,30 (0,26) & $-$39,07\% (19,31\%) & 0,44 (0,39) \\
  MC \textit{every-visit} & 1,17 (0,04) & 0,34 (0,09) & $-$33,24\% (7,94\%)  & 0,47 (0,17) \\
  Q-learning              & 1,31 (0,10) & 0,24 (0,12) & $-$49,58\% (13,24\%) & 0,27 (0,12) \\
  \bottomrule
\end{tabular}

\smallskip
{\footnotesize Média (desvio padrão) em 5 sementes.}
\end{center}

\begin{itemize}
  \item O \textit{every-visit} é o mais estável: menor desvio no teste.
  \item Na semente 42, o \textit{every-visit} venceu no treino (1,20 contra 1,12); na média das 5 sementes, o \textit{first-visit} vence no treino (1,42 contra 1,17). A escolha do \textit{every-visit} dependeu da sorte, embora ele tenha o melhor Sharpe médio e a menor variação no teste.
  \item O mais variável é o MC \textit{first-visit} (desvio de 0,26 no Sharpe do teste), não o Q-learning.
  \item O Q-learning tem o menor Sharpe médio no teste (0,24) e escolhe a mesma ação nas 5 sementes em 21 dos 27 estados.
\end{itemize}

% =============================================================
\section{Resultados no teste}

O teste começa perto de um topo: o BTC caiu de cerca de US\$ 47 mil em janeiro de 2022 para cerca de US\$ 16 mil no fim de 2022. Por isso o \textit{Buy \& Hold} rende pouco, e o foco é o retorno ajustado ao risco.

As políticas são avaliadas no \texttt{gym-trading-env}, um simulador com a interface padrão do Gymnasium: \texttt{env.reset()} devolve a observação inicial, e \texttt{env.step(a)} devolve a observação seguinte, a recompensa e os sinais de fim. A observação é $(\phi(x_t), \phi(x_{t-1}))$, a corretagem é $\kappa$ e o patrimônio começa em US\$ 1.000. O treino do MC e do QL não usa o simulador: percorre o histórico direto, com a mesma recompensa.

\begin{center}
\small
\setlength{\tabcolsep}{5pt}
\begin{tabular}{l cccccc}
  \toprule
  \textbf{Estratégia} & \textbf{Ret. total} & \textbf{Ret. anual} & \textbf{Sharpe} & \textbf{MDD} & \textbf{Calmar} & \textbf{Trocas} \\
  \midrule
  \multicolumn{7}{l}{\textit{Teste (2022 a 2026)}} \\
  MDP (VI)                       & 127,80\%    & 18,37\%    & 0,46    & $-$29,35\% & 0,63    & 113 \\
  Monte Carlo                    & 119,47\%    & 17,47\%    & 0,44    & $-$26,95\% & 0,65    & 108 \\
  Q-learning                     & 86,39\%     & 13,60\%    & 0,26    & $-$51,25\% & 0,27    & 190 \\
  \textit{Buy \& Hold}           & 31,95\%     & 5,84\%     & 0,04    & $-$66,93\% & 0,09    & 1 \\
  Fora do mercado                & 0,00\%      & 0,00\%     & 0,00    & 0,00\%     & --      & 0 \\
  Aleatória                      & $-$93,85\%  & $-$43,52\% & $-$1,45 & $-$94,58\% & $-$0,46 & 1.163 \\
  \midrule
  \multicolumn{7}{l}{\textit{Treino (2014 a 2021)}} \\
  MDP (VI)                       & 6.947,44\%  & 81,24\%    & 1,09    & $-$50,59\% & 1,61    & 140 \\
  Monte Carlo                    & 9.802,30\%  & 90,06\%    & 1,20    & $-$52,95\% & 1,70    & 119 \\
  Q-learning                     & 26.866,84\% & 118,62\%   & 1,37    & $-$69,05\% & 1,72    & 195 \\
  \textit{Buy \& Hold}           & 13.144,05\% & 97,94\%    & 0,90    & $-$83,40\% & 1,17    & 1 \\
  Fora do mercado                & 0,00\%      & 0,00\%     & 0,00    & 0,00\%     & --      & 0 \\
  Aleatória                      & $-$88,73\%  & $-$26,29\% & $-$0,52 & $-$95,74\% & $-$0,27 & 1.721 \\
  \bottomrule
\end{tabular}

\smallskip
{\footnotesize $\kappa = 0{,}1\%$, $\gamma = 0{,}8$; Monte Carlo \textit{every-visit}; Monte Carlo e Q-learning com a semente 42. ``Fora do mercado'' tem Sharpe 0 por convenção: sem volatilidade, o Sharpe não é definido.}
\end{center}

\bloco{Leitura}
\begin{itemize}
  \item \textbf{VI:} maior Sharpe no teste (0,46). \textbf{Monte Carlo:} menor queda máxima entre os RL ($-26{,}95\%$), sem conhecer o modelo.
  \item \textbf{Q-learning:} maior Sharpe no treino (1,37), mas o pior dos três no teste (0,26). Aprendeu a comprar no neutro estável, o que pagou na alta de 2014--2021 e custou caro na queda de 2022. É o problema clássico de generalização: a política funciona no treino e falha fora dele.
  \item Os três RL superam o \textit{Buy \& Hold} em retorno, Sharpe e MDD. VI e MC têm menos da metade da queda máxima do \textit{Buy \& Hold}.
\end{itemize}

\figura{grafico_teste.png}{0.85\textwidth}{Teste: preço do BTC, patrimônio das estratégias e posição escolhida pela política do VI.}

\subsection{Com e sem corretagem}

As mesmas políticas rodam de novo com $\kappa = 0$. As ações são idênticas: a diferença é só o custo das trocas.

\begin{center}
\small
\begin{tabular}{l c cc cc c}
  \toprule
  & & \multicolumn{2}{c}{\textbf{Retorno anual}} & \multicolumn{2}{c}{\textbf{Sharpe}} & \\
  \cmidrule(lr){3-4} \cmidrule(lr){5-6}
  \textbf{Estratégia} & \textbf{Trocas} & \textbf{com} & \textbf{sem} & \textbf{com} & \textbf{sem} & \textbf{Custo (p.p. a.a.)} \\
  \midrule
  MDP (VI)             & 113   & 18,37\%    & 21,46\%    & 0,46    & 0,55    & 3,09  \\
  Monte Carlo          & 108   & 17,47\%    & 20,09\%    & 0,44    & 0,51    & 2,62  \\
  Q-learning           & 190   & 13,60\%    & 18,12\%    & 0,26    & 0,37    & 4,51  \\
  \textit{Buy \& Hold} & 1     & 5,84\%     & 5,87\%     & 0,04    & 0,04    & 0,02  \\
  Aleatória            & 1.163 & $-$43,52\% & $-$21,98\% & $-$1,45 & $-$0,68 & 21,54 \\
  \bottomrule
\end{tabular}
\end{center}

Quem troca mais paga mais: o Q-learning perde 4,51 pontos percentuais por ano. A ordem entre as estratégias não muda.

\subsection{Política gulosa \texorpdfstring{$\times$ $\varepsilon$}{x epsilon}-greedy}

No treino, o robô explora com $\varepsilon$-greedy. Na operação, vale manter $\varepsilon = 0{,}1$ (10\% de ações sorteadas) ou operar de forma gulosa ($a_t = \arg\max_a Q(s_t,a)$)?

\begin{center}
\small
\begin{tabular}{l ccccc}
  \toprule
  \textbf{Política} & \textbf{Sharpe} & \textbf{MDD} & \textbf{Calmar} & \textbf{Trocas} & \textbf{$\varepsilon$ vence a gulosa} \\
  \midrule
  MDP: gulosa                       & 0,46 & $-$29,35\% & 0,63 & 113 & -- \\
  MDP: $\varepsilon$-greedy         & 0,20 & $-$49,87\% & 0,29 & 286 & 24\% das sementes \\
  \midrule
  Monte Carlo: gulosa               & 0,44 & $-$26,95\% & 0,65 & 108 & -- \\
  Monte Carlo: $\varepsilon$-greedy & 0,14 & $-$50,67\% & 0,24 & 289 & 21\% das sementes \\
  \midrule
  Q-learning: gulosa                & 0,26 & $-$51,25\% & 0,27 & 190 & -- \\
  Q-learning: $\varepsilon$-greedy  & 0,04 & $-$59,22\% & 0,13 & 373 & 23\% das sementes \\
  \bottomrule
\end{tabular}

\smallskip
{\footnotesize Teste; $\varepsilon$-greedy com média de 100 sementes.}
\end{center}

A gulosa vence nos três métodos. Com $\varepsilon = 0{,}1$, as trocas mais que dobram (MDP: 113 $\to$ 286) e o Sharpe cai. Na operação, cada ação sorteada é uma decisão que o $Q$ aprendido considera pior, e ainda paga corretagem.

% =============================================================
\section{\textit{Benchmark}: RL \texorpdfstring{$\times$}{x} modelos econométricos}

Seis modelos clássicos, todos estimados só no treino (até 2021). A cada dia, preveem o retorno de amanhã com dados até hoje, e a previsão vira posição: \textbf{previsão de alta $\to$ comprado ($+1$), senão vendido ($-1$)}. Diferente do RL, eles nunca ficam fora do mercado.

\subsection{Passeio aleatório}

\[
  \ln c_{t+1} = \ln c_t + \mu + \varepsilon_{t+1} \quad\Longleftrightarrow\quad \rho_{t+1} = \mu + \varepsilon_{t+1}, \qquad \mathbb{E}[\varepsilon_{t+1}] = 0.
\]
A melhor previsão é $\mu$, estimada pela média do treino. Como $\hat\mu > 0$, a previsão é sempre de alta: comprado todo dia, ou seja, o \textit{Buy \& Hold}. É a hipótese de mercado eficiente na forma fraca: o passado dos preços não ajuda a prever o retorno.

\subsection{AR(1), por mínimos quadrados}

\[
  \rho_{t+1} = \alpha + \phi\,\rho_t + \varepsilon_{t+1}, \qquad
  \hat\phi = \frac{\sum_t (\rho_t - \bar\rho)(\rho_{t+1} - \bar\rho_+)}{\sum_t (\rho_t - \bar\rho)^2}, \qquad \hat\alpha = \bar\rho_+ - \hat\phi\,\bar\rho,
\]
em que $\bar\rho$ e $\bar\rho_+$ são as médias de $\rho_t$ e de $\rho_{t+1}$. $\hat\phi$ é praticamente a autocorrelação de ordem 1 dos retornos.
\begin{itemize}
  \item \textbf{Estacionariedade:} com $|\phi| < 1$, tomando a esperança dos dois lados, a média de longo prazo é $\mathbb{E}[\rho] = \alpha/(1 - \phi)$.
  \item \textbf{Decisão:} comprado se $\hat\alpha + \hat\phi\,\rho_t > 0$. Com $\hat\phi > 0$, segue o movimento de hoje; com $\hat\phi < 0$, aposta na reversão.
\end{itemize}

\subsection{Regressão preditiva (OLS)}

Com $y_t = \rho_{t+1}$ e cada linha de $X$ igual a $(1, x_t, \sigma_{20,t}, \rho_t, \rho_{t-1})$, em que $\sigma_{20,t}$ é a volatilidade de 20 dias:
\[
  y = X\beta + \varepsilon, \qquad \mathrm{SQR}(\beta) = (y - X\beta)'(y - X\beta).
\]

\bloco{Equações normais}
Derivando a soma dos quadrados dos resíduos e igualando a zero:
\[
  \frac{\partial\, \mathrm{SQR}}{\partial \beta} = -2X'(y - X\beta) = 0 \;\Longrightarrow\; X'X\hat\beta = X'y \;\Longrightarrow\; \hat\beta = (X'X)^{-1}X'y.
\]

\bloco{Mínimos quadrados = máxima verossimilhança}
Com $\varepsilon_t \sim N(0, \sigma^2)$ independentes, a log-verossimilhança é
\[
  \ell(\beta, \sigma^2) = -\frac{n}{2}\ln(2\pi\sigma^2) - \frac{\mathrm{SQR}(\beta)}{2\sigma^2},
\]
e maximizar $\ell$ em $\beta$ é o mesmo que minimizar a SQR.

\bloco{Ajuste e significância}
\[
  R^2 = 1 - \frac{\mathrm{SQR}}{\sum_t (y_t - \bar y)^2}, \qquad
  \widehat{\mathrm{Var}}(\hat\beta) = \hat\sigma^2\,(X'X)^{-1}, \qquad
  t_j = \frac{\hat\beta_j}{\mathrm{ep}(\hat\beta_j)}.
\]
No projeto: $R^2 = 0{,}0020$ no treino, e o menor p-valor dos coeficientes é 0,074; nenhum é significativo a 5\%. O modelo explica 0,2\% da variância do retorno de amanhã: o sinal é quase todo ruído, o que ajuda a explicar o pior resultado do \textit{benchmark}.

\subsection{ARMA(1,1), por máxima verossimilhança}

\[
  \rho_{t+1} = \alpha + \phi\,\rho_t + \theta\,\varepsilon_t + \varepsilon_{t+1}, \qquad \varepsilon_t \sim N(0, \sigma^2).
\]
O choque $\varepsilon_t$ não é observado: é reconstruído de forma recursiva, e os parâmetros saem da log-verossimilhança normal, maximizada numericamente:
\[
  \hat\varepsilon_t = \rho_t - \alpha - \phi\,\rho_{t-1} - \theta\,\hat\varepsilon_{t-1}, \qquad
  \ell = -\frac{T}{2}\ln(2\pi\sigma^2) - \frac{1}{2\sigma^2}\sum_t \hat\varepsilon_t^2.
\]
Não há fórmula fechada porque $\hat\varepsilon_t$ depende de $\theta$ de forma recursiva. A previsão é $\hat\rho_{t+1} = \hat\alpha + \hat\phi\,\rho_t + \hat\theta\,\hat\varepsilon_t$; a condição $|\theta| < 1$ faz a recursão ``esquecer'' o valor inicial.

\subsection{Logit e Probit, por máxima verossimilhança}

O alvo é binário, $y_t = 1$ se amanhã sobe, e o modelo dá a probabilidade de alta, $p_t = F(x_t'\beta)$.

\bloco{Logit}
$F$ é a sigmoide, $\sigma(z) = 1/(1 + e^{-z})$. Invertendo, o log das chances é linear: $\ln\big(p_t/(1 - p_t)\big) = x_t'\beta$. Cada dia é uma Bernoulli, $P(y_t) = p_t^{y_t}(1 - p_t)^{1 - y_t}$, então
\[
  \ell(\beta) = \sum_t \big[y_t \ln p_t + (1 - y_t)\ln(1 - p_t)\big],
\]
que é a \textit{binary cross-entropy} com o sinal trocado. Como $\sigma'(z) = \sigma(z)\,(1 - \sigma(z))$:
\[
  \frac{\partial \ell}{\partial \beta} = \sum_t \big[y_t\,(1 - p_t) - (1 - y_t)\,p_t\big]\, x_t = \sum_t (y_t - p_t)\, x_t.
\]
Igualar a zero não tem solução fechada; $\hat\beta$ sai por um método numérico (Newton-Raphson).

\bloco{Probit}
Suponha uma variável não observada $y_t^* = x_t'\beta + u_t$, com $u_t \sim N(0,1)$, e $y_t = 1$ quando $y_t^* > 0$. Pela simetria da normal, $P(y_t = 1) = P(u_t > -x_t'\beta) = \Phi(x_t'\beta)$. A verossimilhança é a mesma, com $\Phi$ no lugar de $\sigma$.

\bloco{Decisão}
Comprado se $\hat p_t > 0{,}5$. Como $\sigma(0) = \Phi(0) = 0{,}5$ e as duas funções são crescentes, isso equivale a $x_t'\hat\beta > 0$. Logit e Probit só diferem pelos $\hat\beta$ estimados, e por isso os resultados quase coincidem.

\subsection{Da previsão ao patrimônio e o resultado}

Todos os modelos, inclusive os RL, são simulados com a mesma fórmula da recompensa, partindo de US\$ 1.000:
\[
  \ell_t = \ln\big(1 + a_t\, r_{t+1} - \kappa\,|a_t - a_{t-1}|\big), \qquad W_T = 1.000 \cdot \exp\Big(\sum_t \ell_t\Big).
\]

\begin{center}
\small
\setlength{\tabcolsep}{5pt}
\begin{tabular}{l ccccc}
  \toprule
  \textbf{Modelo} & \textbf{Ret. total} & \textbf{Sharpe} & \textbf{MDD} & \textbf{Calmar} & \textbf{Trocas} \\
  \midrule
  MDP (\textit{Value Iteration})           & +126,86\% & 0,46    & $-$29,35\% & 0,62    & 113 \\
  Monte Carlo                              & +119,46\% & 0,44    & $-$26,95\% & 0,65    & 108 \\
  Q-learning                               & +86,08\%  & 0,26    & $-$51,28\% & 0,26    & 190 \\
  \midrule
  AR(1)                                    & +110,82\% & 0,23    & $-$73,07\% & 0,23    & 49  \\
  ARMA(1,1)                                & +100,18\% & 0,21    & $-$72,85\% & 0,21    & 43  \\
  Probit                                   & +54,51\%  & 0,10    & $-$72,95\% & 0,13    & 113 \\
  Logit                                    & +50,53\%  & 0,09    & $-$72,95\% & 0,12    & 115 \\
  \textit{Buy \& Hold} (passeio aleatório) & +31,95\%  & 0,04    & $-$66,93\% & 0,09    & 1   \\
  Regressão preditiva (OLS)                & $-$48,48\% & $-$0,35 & $-$77,59\% & $-$0,16 & 225 \\
  \bottomrule
\end{tabular}

\smallskip
{\footnotesize Teste (2022 a 2026), ordenado pelo Sharpe, com a mesma corretagem para todos.}
\end{center}

\bloco{Leitura}
\begin{itemize}
  \item Os três RL lideram em Sharpe e Calmar.
  \item VI e MC têm menos da metade da queda máxima dos econométricos (cerca de $-73\%$).
  \item Os modelos lineares preveem o retorno médio de amanhã sem considerar o custo da troca nem a posição atual. O RL decide com a posição no estado e a corretagem na recompensa, e ficou fora ou vendido em boa parte da queda de 2022.
  \item \textbf{127,80\% ou 126,86\%?} O retorno do VI difere entre a tabela de resultados e o \textit{benchmark} porque são simuladores diferentes: o \texttt{gym-trading-env} não rebalanceia o \textit{short} mantido, e o \textit{benchmark} usa a fórmula exata.
\end{itemize}

% =============================================================
\section{\textit{Walk-forward}: retreinar ao longo do tempo}

O mercado muda de regime. Retreinar o robô de tempos em tempos, só com os dados que já existiam em cada data, deveria melhorar o resultado?

\bloco{Método}
A cada retreino, tudo é refeito: tercis, $\widehat{P}$, $\widehat{R}$, VI, MC, QL e econométricos. Foram 7 configurações: nunca retreinar, ou retreinar a cada 3, 6 ou 12 meses, com todo o histórico (janela expansiva) ou só os últimos 4 anos (janela móvel). Para não escolher olhando o teste, cada modelo fica com a configuração de maior Sharpe numa validação dentro do treino (estimação até 2018, avaliação em 2019--2021).

\bloco{Validação (2019--2021)}
\begin{center}
\small
\begin{tabular}{l c ccc c}
  \toprule
  & & \multicolumn{3}{c}{\textbf{Retreina a cada}} & \\
  \cmidrule(lr){3-5}
  \textbf{Modelo} & \textbf{Nunca} & \textbf{12 meses} & \textbf{6 meses} & \textbf{3 meses} & \textbf{Escolha} \\
  \midrule
  MDP (VI)             & \textbf{1,59} & 0,91          & 0,69 & 1,19          & nunca \\
  Monte Carlo          & 1,58          & 1,22          & 1,46 & \textbf{1,66} & 3 meses \\
  Q-learning           & 1,15          & 0,94          & 0,81 & \textbf{1,17} & 3 meses \\
  AR(1)                & 1,00          & \textbf{1,18} & 1,16 & 1,14          & 12 meses \\
  \textit{Buy \& Hold} & 1,11          & 1,11          & 1,11 & 1,11          & -- \\
  \bottomrule
\end{tabular}

\smallskip
{\footnotesize Sharpe na validação com janela expansiva. A janela móvel não foi escolhida por nenhum desses modelos.}
\end{center}

\bloco{Teste (2022 a 2026)}
\begin{center}
\small
\begin{tabular}{l c ccc c}
  \toprule
  & & \multicolumn{3}{c}{\textbf{Retreina a cada}} & \\
  \cmidrule(lr){3-5}
  \textbf{Modelo} & \textbf{Nunca} & \textbf{12 meses} & \textbf{6 meses} & \textbf{3 meses} & \textbf{Escolha na validação} \\
  \midrule
  MDP (VI)             & 0,46 & 0,28    & 0,35 & 0,07    & nunca \\
  Monte Carlo          & 0,44 & 0,14    & 0,18 & 0,26    & 3 meses \\
  Q-learning           & 0,26 & $-$0,03 & 0,25 & $-$0,17 & 3 meses \\
  AR(1)                & 0,23 & 0,18    & 0,06 & 0,06    & 12 meses \\
  \textit{Buy \& Hold} & 0,04 & 0,04    & 0,04 & 0,04    & -- \\
  \bottomrule
\end{tabular}
\end{center}

\bloco{Leitura}
\begin{itemize}
  \item Com todo o histórico, retreinar \textbf{não melhorou nenhum modelo}.
  \item Exceção: Q-learning só com os últimos 4 anos, retreinando a cada 12 meses (Sharpe 0,33, com uma semente). A validação não escolheu essa configuração.
  \item \textbf{Motivo:} cada retreino recalcula os tercis e as probabilidades. Dias perto da fronteira mudam de faixa, e a política troca de ação em parte dos estados.
  \item Na validação, o ganho de retreinar foi mínimo para o MC (1,58 $\to$ 1,66) e para o QL (1,15 $\to$ 1,17). Diferenças tão pequenas são ruído, e não se repetiram no teste.
\end{itemize}

% =============================================================
\section{Conclusões e limitações}

\bloco{Conclusões}
\begin{itemize}
  \item \textbf{VI:} maior Sharpe no teste (0,46, contra 0,04 do \textit{Buy \& Hold}), com uma política que cabe numa tabela de 27 estados.
  \item \textbf{MC:} chega perto (0,44) sem conhecer o modelo.
  \item \textbf{QL:} mais perto de $Q^*$, mas pior no teste (0,26).
  \item Os três RL lideram em Sharpe e Calmar contra os modelos econométricos.
  \item Retreinar com todo o histórico não ajudou.
\end{itemize}

\bloco{Limitações}
\begin{itemize}
  \item Estado simples: Markov só aproximado (sem volatilidade, volume ou macro).
  \item Um ativo e um único corte entre treino e teste.
  \item Só corretagem: sem \textit{slippage}, impacto de mercado ou aluguel do \textit{short}.
  \item Caixa sem juros no simulador; \textit{short} mantido não rebalanceado no \texttt{gym-trading-env}.
  \item Os tercis e $\widehat{R}$ foram calibrados num período de forte alta (+13.144\% no \textit{Buy \& Hold} do treino): a política herda um viés comprador nos estados estáveis.
  \item Com 5 anos de teste, as diferenças de Sharpe entre os métodos não são estatisticamente significativas.
\end{itemize}

\bloco{Próximos passos}
Volatilidade e volume no estado, outros ativos e custos mais realistas.

\bloco{Lição central}
Com o mesmo estado de 27 estados, os métodos que usam o modelo estimado (VI) ou o retorno real do episódio (MC) generalizaram melhor do que o \textit{bootstrap} do Q-learning. E, no \textit{walk-forward}, a estabilidade da política venceu a tentativa de se adaptar a dados mais recentes.

% =============================================================
\section{Perguntas prováveis e respostas}

\begin{description}[style=nextline, leftmargin=1.2em, labelindent=0pt, itemsep=6pt, font=\bfseries]
  \item[Por que $\gamma = 0{,}8$, se $\gamma = 0{,}3$ ganhou na validação?]
  O 0,8 foi fixado antes, pelo objetivo: a soma dos logs é o log do patrimônio final (critério de Kelly), então o futuro importa. $\gamma = 0{,}3$ é míope (o 5º dia pesa 0,2\%). A bateria só confere a sensibilidade, e entre os $\gamma \ge 0{,}8$ o 0,8 é o melhor.

  \item[Por que não um $\gamma$ baixo, se a recompensa é financeira?]
  É verdade que, com recompensa financeira, a recompensa imediata pode valer mais, e por isso descontamos. Mas $\gamma = 0{,}8$ já desconta forte: o 5º dia pesa 32,8\%. O que se evita é o extremo míope, que ignora que a posição de hoje define o custo de amanhã e que o objetivo é o patrimônio final. E $\gamma < 1$ garante a convergência do VI.

  \item[O mercado é Markov?]
  Não. Os 27 estados são uma aproximação. O MC sofre menos com isso, porque o alvo é o retorno real, e não o valor estimado do próximo estado; mas a política dele também só enxerga os 27 estados. Por isso ele é uma boa comparação com o VI e o QL.

  \item[Por que o Q-learning foi pior, se chega mais perto de $Q^*$?]
  $Q^*$ é ótimo só para o modelo estimado no treino, e pequenas diferenças de valor entre ações já mudam a ação. O QL aprendeu a comprar no neutro estável (premiado na alta do treino, caro em 2022), troca mais (190 contra 113) e depende de Markov ao fazer \textit{bootstrap}. É generalização: bom no treino, ruim fora dele.

  \item[Monte Carlo ou TD: qual tem mais variância?]
  O MC tem alta variância e baixo viés; o TD, menor variância e maior viés. Uma explicação possível para o nosso resultado: o viés do \textit{bootstrap}, somado a um estado só aproximadamente Markov, pesou mais no Q-learning do que a variância pesou no Monte Carlo.

  \item[Qual a diferença entre SARSA e Q-learning?]
  O SARSA usa no alvo a ação realmente tomada (\textit{on-policy}) e aprende o valor da política $\varepsilon$-greedy. O Q-learning usa o máximo (\textit{off-policy}) e aprende o valor da política gulosa, que é a que executamos. O SARSA não foi testado no projeto.

  \item[Por que o passo é $1/N(s,a)$ e não um $\alpha$ fixo?]
  Com $1/N$, $Q(s,a)$ vira a média exata das amostras, e o passo cumpre as condições de convergência. Usar a mesma regra no Monte Carlo e no Q-learning deixa a comparação justa.

  \item[Por que \textit{Value Iteration} e não \textit{Policy Iteration}?]
  Os dois resolvem a mesma equação de otimalidade e chegam ao mesmo $V^*$. Com 27 estados, os dois são rápidos; o VI convergiu em 98 iterações.

  \item[Os resultados dependem da semente?]
  Pouco para o MC \textit{every-visit}: com 5 sementes, o desvio do Sharpe no teste é 0,09. O mais variável é o MC \textit{first-visit} (0,26).

  \item[As diferenças de Sharpe são estatisticamente significativas?]
  Isoladamente, não: com 1.783 dias de teste, o erro padrão do Sharpe anual é cerca de 0,45. O argumento está na consistência: os três RL acima de todos os econométricos em Sharpe e Calmar, com queda máxima menor que a de todos eles, e resultados parecidos com 5 sementes.

  \item[O teste é justo com o \textit{Buy \& Hold}?]
  Todos operam no mesmo período, que começa perto de um topo. Por isso o foco é o retorno ajustado ao risco: os RL vencem em Sharpe e em queda máxima.

  \item[Por que o \textit{Buy \& Hold} aparece como passeio aleatório?]
  O passeio aleatório tem média positiva, então prevê alta todo dia e fica sempre comprado: é exatamente o \textit{Buy \& Hold}.

  \item[Por que o retorno do VI é 127,80\% numa tabela e 126,86\% na outra?]
  Simuladores diferentes: o \texttt{gym-trading-env} não rebalanceia o \textit{short} mantido; o \textit{benchmark} usa a fórmula exata para todos.

  \item[Por que retreinar não ajudou?]
  Cada retreino muda tercis e probabilidades, e a política troca de ação. Os ganhos na validação foram mínimos (MC +0,08, QL +0,02) e não se repetiram no teste.

  \item[O robô aprende online ou offline?]
  Offline: aprende com o histórico de preços, sem operar de verdade. Funciona porque o robô não mexe no preço: qualquer ação pode ser simulada sobre o histórico, e o que depende da ação (posição e corretagem) é calculado exatamente. A desvantagem, ser menos adaptável, é o que o \textit{walk-forward} testa.

  \item[Por que a recompensa em log?]
  A soma dos logs diários é o log do patrimônio final (aditividade exata), e a recompensa não depende do tamanho da conta.

  \item[Por que 365 dias e não 252?]
  O Bitcoin negocia todos os dias; cada dia é um período.

  \item[Por que tercis e não limites fixos?]
  Dividem o treino em faixas com 1/3 dos dias cada, são calculados só no treino (sem vazar o teste) e se adaptam à distribuição do ativo.

  \item[Por que o estado $(0,2)$ mantém a posição?]
  Ele nunca ocorre no treino: $\widehat{P}$ é uniforme e o retorno é nulo, então a única recompensa é o custo da troca.

  \item[Por que a $\varepsilon$-greedy perde?]
  Com 10\% de ações sorteadas, as trocas mais que dobram e o robô fica em posições erradas. A versão com exploração venceu a gulosa em só 21\% a 24\% das sementes.
\end{description}

\end{document}
